%! Addition and Multiplication Rules — Unit 3, Lesson 3.3 — Student Packet Cover Sheet
% THREE packet rows — Warm-Up, Guided Notes & Practice, Homework. The homework is the individual
% practice, assigned at Close & Assign and done outside class, and it is scored 1 to 5.
\documentclass[12pt]{article}
\usepackage{saar-article}
\usepackage{saar-boxes}
\usepackage{ltablex}
\keepXColumns
% Measured banner: the band grows to fit the title block, so a title long enough to wrap grows
% the band rather than running out of it (the stats course's \coverbanner; saar-article does not
% carry one, so it is defined here — every cover copies this block verbatim).
\newsavebox{\bannerbox}
\newlength{\bannerht}
\newlength{\coverbannerpad}
\setlength{\coverbannerpad}{0.16in}       % breathing room above and below the text
\newcommand{\coverbanner}[2]{%
  \sbox{\bannerbox}{%
    \begin{minipage}{\dimexpr\paperwidth-1.4in\relax}%
      \centering
      {\color{white}\textbf{\LARGE \CourseName}}\\[4pt]
      {\color{frostmid}\large\textbf{#1}}\\[6pt]
      {\color{white}\textbf{\large #2}}%
    \end{minipage}}%
  \setlength{\bannerht}{\dimexpr\ht\bannerbox+\dp\bannerbox+2\coverbannerpad\relax}%
  \noindent\begin{tikzpicture}[remember picture, overlay]
    \fill[cerulean] (current page.north west)
      rectangle ([yshift=-\bannerht]current page.north east);
    \node[anchor=north, inner sep=0pt]
      at ([yshift=-\coverbannerpad]current page.north) {\usebox{\bannerbox}};
  \end{tikzpicture}\par
  % The text body starts 0.6in below the page edge (geometry top=0.6in), so reserve only what
  % the band overruns that, plus a gap under the band.
  \vspace*{\dimexpr\bannerht-0.6in+0.16in\relax}%
}

\begin{document}

% Full-bleed cerulean banner, sized to the title block.
\coverbanner{Unit 3}{Lesson 3.3 \quad Addition and Multiplication Rules}

\vspace{0.06in}
% NAMESTRIP: this is the ONLY component that carries the name row.
\namedateperiod
\vspace{0.18in}

% The standard codes (PS.P.1a, c; AFDA.DA.3e, f) stay in the teacher-facing plan.
\begin{learningtargetbox}
\small
\begin{enumerate}[leftmargin=1.5em, itemsep=5pt]
  \item \ldots{} use the \textbf{addition rule} on a Venn diagram, subtracting the overlap, and
        drop the subtraction when events are \textbf{mutually exclusive}.
  \item \ldots{} use the complement rule to find the probability of \emph{at least one}.
  \item \ldots{} use the \textbf{multiplication rule} on a probability tree for
        \textbf{independent events} (with replacement) and \textbf{dependent events} (without
        replacement).
  \item \ldots{} decide by computing whether two events are complementary, mutually exclusive,
        independent, or dependent --- and explain why mutually exclusive events are not
        independent.
\end{enumerate}
\end{learningtargetbox}

\vspace{0.14in}

\begin{tocbox}
\small
\renewcommand{\arraystretch}{1.5}
\begin{tabularx}{\linewidth}{@{} c l X @{}}
\rowcolor{frost}
\textbf{\#} & \textbf{Component} & \textbf{Description} \\
\midrule
1 & Warm-Up & Adding and multiplying decimals and fractions, then reading the concert raffle's
             Venn diagram \\
2 & Guided Notes \& Practice & The concert raffle: the addition rule, mutually exclusive
             events, the multiplication rule on trees, and classifying pairs of events --- then
             rain and a late bus, worked together \\
% Row 3 is the lesson's GRADED practice, assigned at the close. The row never changes with the form.
3 & Homework & The prize wheel and capsule tray at an arcade: the rules, then which pairs of
             events are independent. \textbf{Scored 1 to 5, assigned today, due the first
             class after two study halls.} \\
\end{tabularx}
\end{tocbox}

\vspace{0.14in}

% Keep in Mind carries the lesson's CONTENT — the rule, the test, the trap — never its process.
\begin{remindbox}
\small
\textbf{Addition rule:} $P(A \text{ or } B) = P(A) + P(B) - P(A \text{ and } B)$; if the events
are \textbf{mutually exclusive}, the overlap is $0$. \textbf{At least one:}
$1 - P(\text{none})$. \textbf{Multiplication rule:} multiply along the branches of a tree; the
second factor stays the same for \textbf{independent events} (with replacement) and changes for
\textbf{dependent events} (without replacement). \textbf{Watch out:} events are independent when
$P(A \text{ and } B) = P(A) \cdot P(B)$. Mutually exclusive events that each have some chance
are \textbf{dependent}: knowing one happened tells you the other did not.
\end{remindbox}

\end{document}
